%!TeX root=MemoriaTFG.tex

\chapter{Limitaciones}
\label{cap:limitaciones}

El estudio presenta limitaciones de cinco órdenes que conviene
declarar explícitamente. La presente sección las desarrolla por
familia, cuantifica su magnitud cuando es posible y acota su
impacto sobre la interpretación de los resultados del
capítulo~\ref{cap:resultados} y de la discusión del
capítulo~\ref{cap:discusion}.

\section{Limitaciones de disponibilidad}
\label{sec:lim-disponibilidad}

\subsection{Modelos sin pesos públicos}
\label{sec:lim-pesos}

Los pesos de DermFM-Zero~\cite{dermfmzero2026} no son públicos en
la fecha de cierre del presente trabajo. Esta restricción impide
su evaluación empírica directa sobre los doce datasets externos
contrastados y restringe su tratamiento al marco conceptual del
capítulo~\ref{cap:sota}. La comparación entre PanDerm Large y
DermFM-Zero, que la publicación original reporta a favor del
segundo, no puede verificarse de forma independiente con los mismos
protocolos. La disponibilidad eventual de los pesos permitiría esa
comparación directa con los hallazgos cuantitativos del trabajo
(tabla~\ref{tab:lp-panderm}, tabla~\ref{tab:equidad-fototipo}).

Los modelos cerrados accedidos vía API (GPT-4o, GPT-4o-mini,
Gemini~2.5~Pro, Gemini~2.5~Flash) presentan una limitación
estructural: no permiten controlar la versión exacta del modelo en
el momento de la inferencia ni el posprocesado interno aplicado
por el proveedor. Las cifras de la tabla~\ref{tab:llm-comparativa}
son condicionales a la versión disponible en la fecha del
experimento (marzo de 2026) y no son estrictamente reproducibles
ante actualizaciones silenciosas del proveedor. El coste agregado
de la evaluación (\$4{,}07 aproximadamente entre GPT-4o y
GPT-4o-mini, sin desglose público de tarifa por token de entrada y
salida) introduce una segunda dependencia que limita la
replicación por terceros.

\subsection{Conjuntos de preentrenamiento no liberados}
\label{sec:lim-preentrenamiento}

Los conjuntos de preentrenamiento internos del grupo de la Universidad
Monash empleados en PanDerm~\cite{panderm2025} no son accesibles para
terceros, lo que impide reproducir la fase de preentrenamiento y limita
la verificación a los pesos liberados. En consecuencia, cualquier
afirmación sobre la composición del dataset de preentrenamiento (las
cuatro modalidades dermatológicas, $\sim 2{,}1$ millones de imágenes)
descansa sobre la declaración de los autores y no sobre auditoría
independiente. La misma restricción aplica al dataset Derm1M de la
familia DermLIP, distribuido con \emph{embeddings} agregados pero sin
acceso íntegro al material de imagen original.

\subsection{Dependencia de modelos fundacionales externos}
\label{sec:lim-dependencia-externa}

El trabajo evalúa principalmente estrategias de adaptación,
transferencia e integración sobre modelos fundacionales desarrollados
por terceros (PanDerm, DermLIP, SAM2.1, SigLIP, BiomedCLIP). Por tanto,
parte de las conclusiones depende de las decisiones de diseño y de los
sesgos heredados de dichos modelos. Aunque esta situación refleja el
escenario real de adopción clínica actual ---en el que una institución
sanitaria parte de pesos preentrenados disponibles y no de un
preentrenamiento propio a gran escala---, limita la capacidad de atribuir
causalmente determinados comportamientos a componentes específicos del
entrenamiento original.

\section{Limitaciones de los datos}
\label{sec:lim-datos}

\subsection{Sesgo geográfico y lingüístico}
\label{sec:lim-geografico}

Los doce datasets externos contrastados provienen mayoritariamente
de instituciones de Estados Unidos (DDI), Australia
(HAM10000 parcial), Europa central (BCN20000 sobre Hospital Clínic
de Barcelona; HAM10000 parcial; ISIC2018; Dermnet) y Brasil
(PAD-UFES-20). HIBA procede del Hospital Italiano de Buenos Aires
y MSKCC del Memorial Sloan Kettering Cancer Center. La población
hispanohablante europea, con la distribución epidemiológica
documentada en DIADERM~\cite{diaderm2018}, no está representada de
forma específica en los conjuntos de evaluación. La consecuencia
es que las cifras agregadas de la tabla~\ref{tab:lp-panderm} no
son directamente extrapolables al espectro diagnóstico atendido en
el Hospital Universitario Son Llàtzer
(HUSLL)~\cite{taberner2010_motivos} sin validación adicional.

La rama lingüística de los modelos vision-lenguaje contrastivos
(DermLIP, BiomedCLIP) está preentrenada sobre dataset en inglés. La
sensibilidad documentada al tokenizador (sección~\ref{sec:disc-zs},
$48{,}8$\,pp AUROC entre la peor y la mejor configuración) se
limita a las variantes evaluadas en inglés y no caracteriza el
rendimiento sobre \emph{prompts} en español, idioma sobre el que
no se han evaluado por la ausencia de validación de los
tokenizadores con vocabulario clínico hispanohablante. El dataset
DermapixelAI~1.0 publicado como contribución del trabajo
(sección~\ref{sec:dermapixel}) aborda esta brecha para la fase
posterior al periodo cubierto por este documento.

\subsection{Heterogeneidad taxonómica entre datasets}
\label{sec:lim-taxonomia}

Las taxonomías diagnósticas nativas de los doce datasets son
incompatibles entre sí. La formulación binaria
melanoma/no-melanoma de Derm7pt clínico, Derm7pt dermo, HIBA y
MSKCC contrasta con la multiclase de HAM10000 ($7$ clases),
PAD-UFES-20 ($6$), DDI ($2$), Dermnet ($23$), BCN20000
($9$), WSI~patches ($16$) y Fitzpatrick17k
($114$ patologías). SkinCon emplea anotación basada en
$48$ conceptos clínicos visuales, taxonomía incompatible con la
diagnóstica de los demás datasets.

La ontología jerárquica L1/L2/L3 introducida en
sección~\ref{sec:ontologia} mitiga esta limitación al armonizar el
dataset agregado sobre vocabulario común
($72\,654$ imágenes etiquetadas, $4$ categorías L1, $43$
subcategorías L2 con cinco entradas sin desagregación a L3, y
$367$ diagnósticos L3 con codificación SNOMED~CT y CIE-10), pero
no la elimina por completo. Las comparaciones inter-dataset son
defendibles a nivel L1 (cuatro familias etiológicas) y L2
($43$ entradas), pero a nivel L3 cada dataset conserva su cola
larga propia y los contrastes directos requieren agregación
adicional. Once de los doce datasets se mapean a la ontología;
SkinCon queda fuera por el desajuste de granularidad ya
mencionado.

La consecuencia operativa para los hallazgos del
capítulo~\ref{cap:resultados} es que el reporte por dataset de la
tabla~\ref{tab:lp-panderm} mantiene fidelidad a la taxonomía
nativa, y las cifras agregadas (filas \emph{Media} de la misma
tabla) deben interpretarse como promedios sobre tareas heterogéneas
no estrictamente comparables entre sí.

\subsection{Solapamiento con el dataset de preentrenamiento Derm1M}
\label{sec:lim-leakage}

La auditoría descrita en sección~\ref{sec:auditoria} identifica un
único dataset evaluado, \textbf{Dermnet}, con solapamiento exacto
frente al dataset de preentrenamiento Derm1M empleado por la familia
DermLIP: sus $18\,302$ imágenes están íntegramente contenidas en
Derm1M ($100\%$ por hash MD5). HIBA y MSKCC, pese a su origen en el
ISIC archive, no presentan ninguna coincidencia MD5 con Derm1M, por
lo que la sospecha de procedencia no se confirma como solapamiento
exacto (tabla~\ref{tab:datasets-resumen}).

Las métricas obtenidas sobre Dermnet con modelos cuya rama visual
deriva de Derm1M (DermLIP v1, v2 y original) deben interpretarse con
\emph{caveat} y no son directamente comparables con las obtenidas
sobre los datasets sin solapamiento. La sensibilidad de la métrica
AUROC a la presencia de imágenes ya vistas en preentrenamiento puede
inflar al alza las cifras zero-shot de DermLIP sobre Dermnet.

Las cifras de PanDerm Base y Large, DINOv2 ViT-L/14,
ConvNeXt-Large, EfficientNetV2-Large y SigLIP-Large SO400M no se
ven afectadas por este solapamiento, dado que su preentrenamiento
no incluye Derm1M. Por consiguiente, los hallazgos cuantitativos
sobre la ventaja del preentrenamiento de dominio
(sección~\ref{sec:disc-especializacion}), la eficiencia en etiquetas
(sección~\ref{sec:disc-eficiencia}) y la comparación con CNNs
(tabla~\ref{tab:gap-cnn}) permanecen defendibles incluso bajo el
peor escenario de la auditoría.

La auditoría acota, no obstante, su propio alcance epistemológico. El
hash MD5 detecta coincidencias \emph{exactas} a nivel de bytes, pero
no identifica \emph{near-duplicates}: recortes, reescalados,
recompresiones o alteraciones leves de una misma imagen. Lo que la
auditoría sostiene es, por tanto, la ausencia de duplicación
\emph{exacta} entre los datasets de evaluación y Derm1M, no la
ausencia completa de contaminación visual. Las cifras declaradas
libres de solapamiento deben leerse en sentido estricto como libres de
\emph{solapamiento exacto}; la extensión natural ---\emph{perceptual
hashing} o comparación por \emph{embeddings} visuales--- se identifica
como refinamiento metodológico de la auditoría.

\subsection{Tamaños muestrales por subgrupo en el experimento de equidad}
\label{sec:lim-muestrales}

Los tamaños muestrales del experimento de equidad sobre
Fitzpatrick17k son asimétricos entre fototipos cutáneos: $299$
imágenes para FST~I, $470$ para FST~II, $325$ para FST~III,
$261$ para FST~IV, $169$ para FST~V y $73$ para FST~VI
en el conjunto de test ($N$ total $=1\,597$ con fototipo
asignado). El subgrupo FST~VI dispone de aproximadamente la cuarta
parte de los ejemplos del subgrupo FST~I y supone el $4{,}6\%$
del conjunto de test.

La consecuencia estadística es que los intervalos de confianza al
$95\%$ por \emph{bootstrap} sobre las cifras de
BAcc por subgrupo de la tabla~\ref{tab:equidad-fototipo} son más
amplios para FST~V-VI que para FST~I-IV, lo que limita la
precisión cuantitativa del \emph{gap}~I-VI reportado. La
interpretación cualitativa ---rendimiento inferior sobre fototipos
elevados--- es robusta frente a esta asimetría muestral; la
cuantificación exacta del \emph{gap}, en cambio, queda condicionada
por la menor masa muestral de los fototipos oscuros, y su
acompañamiento sistemático con intervalos de confianza por
\emph{bootstrap} se identifica como refinamiento estadístico de la
línea de trabajo.

\section{Limitaciones del diseño experimental}
\label{sec:lim-diseno}

\subsection{Variabilidad inter-semilla}
\label{sec:lim-semilla}

Los protocolos de sondeo lineal, eficiencia en etiquetas y
evaluación de equidad por fototipo operan con una única semilla
(\texttt{random\_state}~$=42$); el protocolo de ajuste fino emplea
también una única configuración de
semillas (\texttt{torch}, \texttt{numpy}, \texttt{random} fijadas a
$0$). La variabilidad inter-semilla no se reporta de forma
sistemática en estos protocolos, lo que limita la cuantificación de
la incertidumbre estadística asociada a la inicialización aleatoria
del clasificador lineal y al orden de muestreo durante el
entrenamiento. La segmentación constituye la excepción: el
cifra principal de Dice se acompaña de una validación cruzada de
cinco pliegues con su intervalo de confianza
(sección~\ref{sec:seg-robustez}), por lo que en esa tarea la
incertidumbre de partición sí queda acotada.

La consecuencia es que las cifras de las tablas
\ref{tab:lp-panderm}, \ref{tab:ft-panderm}, \ref{tab:label-eff} y
\ref{tab:equidad-fototipo} son estimaciones puntuales sobre una
trayectoria de optimización concreta y no incluyen su variabilidad
esperable bajo distintas inicializaciones. Las heurísticas
operativas extraídas en sección~\ref{sec:disc-lp-vs-ft}
(régimen $N_{\mathrm{train}} \lesssim 500$ para preferencia de
sondeo lineal) son cualitativas y robustas, pero los umbrales
exactos pueden desplazarse bajo otras semillas.

Esta limitación se acota por dos consideraciones. La semilla fija
respondió al objetivo de comparar arquitecturas bajo un protocolo
constante, aislando el efecto del modelo del de la inicialización. Y
la magnitud de los efectos cualitativos sobre los que se sostienen las
conclusiones ---del orden de $+9{,}0$\,pp en la ventaja del
preentrenamiento de dominio o de $43{,}5$\,pp en la brecha entre
encoders especializados y LLMs generalistas--- excede ampliamente la
variabilidad inter-semilla esperable para un clasificador lineal sobre
estos tamaños muestrales, por lo que no se anticipan cambios
cualitativos en las conclusiones. El reporte sistemático de media y
desviación estándar sobre múltiples semillas constituye, en todo caso,
la extensión natural de este protocolo.

\subsection{Pruebas estadísticas y comparaciones múltiples}
\label{sec:lim-estadistica}

El test de DeLong~\cite{delong1988} se ha aplicado de forma sistemática
a los contrastes pareados entre los cinco encoders sobre el
\emph{endpoint} binario de detección de malignidad en Fitzpatrick17k,
con corrección de Bonferroni para las diez comparaciones
(sección~\ref{sec:equidad-delong}, tabla~\ref{tab:delong-pares}). Su
extensión al resto de \emph{endpoints} y tablas del capítulo, junto con
la presentación de cada cifra principal como par (estimador,
IC~\emph{bootstrap}) en lugar de valor puntual, permanece como
refinamiento estadístico de la versión publicable.

La consecuencia operativa es que la superioridad entre encoders sobre la
detección de malignidad queda establecida estadísticamente, mientras que
el resto de hallazgos numéricos del capítulo~\ref{cap:resultados} deben
interpretarse en términos descriptivos y de magnitud del efecto: las
afirmaciones de \emph{superioridad estadísticamente significativa} en los
demás \emph{endpoints} requieren la cuantificación pendiente. Los
patrones cualitativos identificados en la discusión son, no obstante,
robustos frente a esta limitación por la magnitud absoluta de los efectos
cuantificados.

\subsection{Sensibilidad al \emph{prompt} en zero-shot y LLMs}
\label{sec:lim-prompts}

La evaluación zero-shot multimodal y la comparación con LLMs
multimodales operan sobre formulaciones concretas de \emph{prompts}
que condicionan los resultados reportados. La diferencia de
$48{,}8$\,pp AUROC entre la peor y la mejor configuración de
DermLIP sobre HAM10000 (tabla~\ref{tab:zs-ham}) cuantifica la
magnitud de esta dependencia para la rama contrastiva. La
generalización a otras formulaciones requiere validación
específica.

Para los LLMs multimodales, el \emph{prompt} clínico de
aproximadamente $300$ palabras descrito en sección~\ref{sec:llm}
introduce una formulación fija que no se ha auditado mediante
barrido sistemático de longitud, orden de las clases candidatas o
inclusión de descriptores anatómicos. El sesgo léxico documentado
en sección~\ref{sec:disc-llm} (GPT-4o emite $284$ predicciones de
dermatofibroma frente a $8$ muestras reales) sugiere que el
modelo responde al espacio léxico del \emph{prompt}, lo que
amplifica el efecto de la formulación. Las cifras de la
tabla~\ref{tab:llm-comparativa} son condicionales a este
\emph{prompt} y no caracterizan el rendimiento de los modelos bajo
formulaciones alternativas. Una caracterización sistemática de la
sensibilidad al \emph{prompt} en LLMs queda fuera del alcance del
trabajo y se documenta como línea futura en el
Anexo~\ref{cap:anexog}.

\section{Limitaciones del dataset DermapixelAI 1.0}
\label{sec:lim-dermapixel}

El dataset DermapixelAI~1.0 entregado como contribución del trabajo
presenta cuatro limitaciones específicas declaradas explícitamente
en el Anexo~\ref{cap:anexoc} y en el documento formal de entrega
(Anexo~\ref{cap:anexoe}).

\paragraph{Desbalance estructural por modalidad.}
La distribución por modalidad es $1\,036$ imágenes de fotografía
clínica frente a $49$ imágenes dermatoscópicas, lo que limita
el uso del dataset para entrenamiento o evaluación de tareas
dermatoscópicas. El predominio clínico refleja la finalidad
docente original del archivo Dermapixel, no una decisión de
diseño del trabajo, pero condiciona la utilidad del dataset para
benchmarks comparativos con HAM10000, BCN20000, ISIC2018 u otros
conjuntos dermatoscópicos.

\paragraph{Cobertura ontológica parcial.}
El dataset contiene anotación L3 sobre $250$ diagnósticos
distintos, sobre las $367$ entradas declaradas en la
ontología, lo que corresponde a una cobertura efectiva del
$68{,}1\%$ del vocabulario L3 con al menos una imagen. El
diagnóstico L3 más frecuente concentra aproximadamente el
$32\%$ del dataset, lo que reproduce la cola larga típica del
material clínico. Esta cobertura parcial limita el uso del dataset
para evaluación granular sobre L3 sin agregación a L2.

\paragraph{Tres cifras de validación no intercambiables.}
La distinción entre los tres campos de validación documentados en
el Anexo~\ref{cap:anexoc} es relevante para la interpretación
correcta de las cifras de calidad. El $97{,}89\%$ del dataset
con \texttt{label\_source = ontology} indica mapeo ontológico
validado del vocabulario, no verificación visual caso-a-caso. El
$84{,}39\%$ con \texttt{rosa\_verified} corresponde a la revisión
visual explícita por la Dra.~R.~Taberner. La cifra de validación
textual \texttt{diagnosis\_source = expert\_v3} ($98{,}38\%$) se
refiere a la coherencia del razonamiento textual con el
diagnóstico asignado, sin verificación visual de la imagen
asociada. Las tres cifras deben declararse con precisión en
cualquier comunicación derivada y no son intercambiables.

\paragraph{Splits y tamaños reducidos.}
La partición \emph{case-aware} del dataset
($891 / 157 / 41$ imágenes en \emph{train/val/test}) es
metodológicamente correcta para evitar fuga de información entre
imágenes del mismo caso clínico, pero produce un conjunto de test
reducido ($41$ imágenes en el \emph{split} de producción; los
contrastes por nivel ontológico se evalúan sobre el subconjunto
filtrado, $N = 36$ en L1/L2 y $N = 28$ en L3) que limita la potencia
estadística de los contrastes realizados sobre él. La explotación experimental
sistemática sobre DermapixelAI~1.0 se ha pospuesto al periodo
posterior al cierre de este documento por esta razón. En consecuencia,
las conclusiones comparativas sobre el dataset ---en particular la
jerarquía entre métodos en L2 y L3 y la narrativa de Dermapixel R0---
se anclan en la validación cruzada agrupada por caso
(sección~\ref{sec:disc-cv-robustez}, tabla~\ref{tab:dermapixel-cv}) y
no en la cifra puntual del \emph{split} fijo, que sobrestima la BAcc
sobre este \emph{test} reducido.

\section{Limitaciones de generalización clínica}
\label{sec:lim-clinica}

\subsection{Material de archivo frente a flujo asistencial real}
\label{sec:lim-archivo}

El estudio mide rendimiento sobre fotografías procedentes de
archivos de investigación curados (HAM10000, BCN20000, ISIC2018) o
de archivos clínicos con verificación posterior (HIBA, MSKCC, DDI,
Fitzpatrick17k). El rendimiento sobre material adquirido en el
flujo real de teledermatología, sobre fotografías de atención
primaria con dispositivos heterogéneos (variabilidad de
iluminación, distancia focal, encuadre, calibración cromática) o
sobre material adquirido en consulta presencial sin estandarización
no se aborda en este trabajo. La extrapolación de las cifras de la
tabla~\ref{tab:lp-panderm} a estos escenarios requiere validación
adicional sobre cohorte representativa del entorno de despliegue
previsto, particularmente en el contexto del proyecto institucional
de teledermatología con IB-Salut descrito en
sección~\ref{sec:trabajo-previo}.

\subsection{Sesgo de modalidad de los componentes evaluados}
\label{sec:lim-modalidad-sistema}

Varios de los protocolos evaluados se apoyan en dataset
predominantemente dermatoscópicos: la clasificación adaptada sobre
HAM10000, la segmentación promptable adaptada sobre ISIC2018
(sección~\ref{sec:seg-results}) y la lista de comprobación de
siete puntos, cuyos criterios son intrínsecamente dermatoscópicos.
Dado que la entrada mayoritaria en un escenario de teledermatología
real es la fotografía clínica ---como refleja la propia composición
del dataset DermapixelAI~1.0, con $1\,036$ imágenes clínicas frente a
$49$ dermatoscópicas (sección~\ref{sec:lim-dermapixel})---, la
transferencia de estos componentes a fotografía clínica no
estandarizada queda como limitación operativa conocida. Su
mitigación ---clasificación previa de la modalidad de imagen y
encaminamiento diferenciado por tipo--- se identifica como línea de
continuación y no se aborda en este trabajo.

\subsection{Ausencia de validación prospectiva}
\label{sec:lim-prospectivo}

El trabajo no aborda la validación clínica prospectiva sobre
cohorte hospitalaria real con protocolo aprobado por Comité Ético
de Investigación Clínica (CEIC). Esta limitación es estructural y
se asume explícitamente como decisión de alcance en
sección~\ref{sec:alcance}. La validación prospectiva del prototipo
DermApIxel descrito en el Anexo~\ref{cap:anexoh} sobre cohorte del
HUSLL requiere prerrequisitos administrativos y temporales
(aprobación CEIC, definición de \emph{endpoints} clínicos
primarios y secundarios, cálculo muestral, contratación de
revisores independientes) que exceden el alcance del Trabajo de
Fin de Grado y se identifican como línea de continuación inmediata
en la sección~\ref{sec:lineas-abiertas} del capítulo~\ref{cap:conclusiones}.

\subsection{Elementos no abordados por decisión de alcance}
\label{sec:lim-fuera}

Cuatro elementos quedan fuera del alcance declarados en
sección~\ref{sec:alcance} y deben reconocerse como limitaciones del
documento aunque sean por decisión deliberada: (i)~el
preentrenamiento desde cero de un modelo fundacional propio, que
requeriría acceso al material de las cuatro modalidades
dermatológicas en volumen comparable al de PanDerm
($\sim 2{,}1$ millones de imágenes); (ii)~la integración
operativa con el sistema PACS y el visor radiológico del HUSLL,
que excede el alcance técnico del trabajo experimental;
(iii)~el reentrenamiento multilingüe de la rama lingüística sobre
dataset en español de gran escala, condicionado a la disponibilidad
futura de material en cantidad suficiente; (iv)~la certificación
CE como producto sanitario (\emph{Software as a Medical Device}),
que requiere un proceso regulatorio y de auditoría independiente
no contemplado.

\section{Síntesis: dominio de validez de los resultados}
\label{sec:lim-sintesis}

El conjunto de limitaciones declarado en este capítulo acota el
dominio de validez de las cifras reportadas en el
capítulo~\ref{cap:resultados} pero no invalida los patrones
cualitativos identificados en la discusión. Cinco afirmaciones resisten
al peor escenario combinado de muestreo, semilla, sensibilidad al
\emph{prompt} y restricciones de modelos cerrados, y constituyen la
respuesta empírica defendible del trabajo a la pregunta de investigación
formulada en sección~\ref{sec:pregunta}:

\begin{enumerate}
  \item La \textbf{ventaja del preentrenamiento sobre el dominio
  clínico} frente a codificadores generalistas.
  \item La \textbf{saturación temprana del sondeo lineal}, que alcanza
  el grueso de su rendimiento con $1$--$5\%$ de los datos etiquetados.
  \item La \textbf{brecha entre codificadores específicos y LLMs
  multimodales generalistas}, a favor de los primeros.
  \item El \textbf{\emph{gap} sistemático por fototipo cutáneo}, con
  sentido peor-en-fototipos-elevados y consistente entre encoders.
  \item La \textbf{ausencia de un método óptimo único y la consiguiente
  especialización por tarea} (sección~\ref{sec:disc-especializacion}),
  lectura transversal que sintetiza las cuatro anteriores.
\end{enumerate}

Las cuatro primeras son afirmaciones directas; la quinta es una lectura
de orden superior que emerge de leerlas en conjunto. Su cuantificación y
su lectura como conclusiones se desarrollan en el
capítulo~\ref{cap:conclusiones} (sección~\ref{sec:respuesta}). La quinta
está, además, acotada por el tamaño moderado del dataset
DermapixelAI~1.0 (sección~\ref{sec:lim-dermapixel}), su cola larga en el
nivel L3 y el carácter exploratorio de varios de los contrastes por
nivel ontológico: el patrón cualitativo es robusto, pero no constituye
una ordenación cuantitativa cerrada entre métodos, cuya validación
requeriría los tamaños muestrales y la cohorte prospectiva descritos en
el capítulo~\ref{cap:conclusiones}.
